\documentclass[11pt,a4paper]{article}
% ============================================================
%  פרמבולה משותפת לחומרי הבחינה — מודולים 16–18
%  מהדרים עם:  lualatex <file>.tex
% ============================================================
\usepackage{fontspec}
\directlua{
  luaotfload.add_fallback("hbfb", {"DejaVu Serif:mode=node;", "DejaVu Sans:mode=node;"})
  luaotfload.add_fallback("ltfb", {"DejaVu Sans:mode=node;"})
}
\usepackage[bidi=basic]{babel}
\babelprovide[import, main]{hebrew}
\babelprovide[import]{english}
\babelfont{rm}[Ligatures=TeX, RawFeature={fallback=ltfb}]{DejaVu Serif}
\babelfont[hebrew]{rm}[Script=Hebrew, Ligatures=TeX, RawFeature={fallback=hbfb}]{Noto Serif Hebrew}
\babelfont{tt}[Scale=0.84, RawFeature={fallback=ltfb}]{DejaVu Sans Mono}

\usepackage[a4paper,top=2.1cm,bottom=2.0cm,left=1.9cm,right=1.9cm,headsep=0.5cm,footskip=1cm]{geometry}
\usepackage{booktabs, longtable, array, adjustbox, xcolor, enumitem, fancyhdr, titlesec,
            fancyvrb, needspace, multicol, lastpage, tikz, colortbl}
\usetikzlibrary{positioning, fit, backgrounds, matrix}

\definecolor{inkblue}{HTML}{14365A}
\definecolor{rulegray}{HTML}{8A97A3}
\definecolor{boxgray}{HTML}{F4F6F8}
\definecolor{codebg}{HTML}{F7F8FA}
\definecolor{warnred}{HTML}{9B2226}
\definecolor{okgreen}{HTML}{1B5E3A}
\definecolor{tipgold}{HTML}{8A6D1F}

\setlength{\parindent}{0pt}
\setlength{\parskip}{3pt plus 1pt}
\renewcommand{\arraystretch}{1.15}
\setlength{\columnsep}{18pt}
\setlength{\headheight}{13pt}

\pagestyle{fancy}\fancyhf{}
\renewcommand{\headrulewidth}{0.4pt}
\renewcommand{\footrulewidth}{0pt}
\fancyfoot[C]{\small עמוד \thepage{} מתוך \pageref{LastPage}}

% ---- כותרות ----
\titleformat{\section}{\needspace{3\baselineskip}\large\bfseries\color{inkblue}}{\thesection.\ }{0pt}{}
\titlespacing*{\section}{0pt}{15pt plus 3pt}{6pt}
\titleformat{\subsection}{\needspace{2.5\baselineskip}\normalsize\bfseries\color{inkblue}}{\thesubsection\ }{0pt}{}
\titlespacing*{\subsection}{0pt}{10pt plus 2pt}{4pt}

% ---- קו מפריד מתחת לכותרת חלק ----
\newcommand{\partheader}[3]{%
  \section*{#1}%
  \vspace{-4pt}{\color{rulegray}\hrule height 0.8pt}\vspace{4pt}%
  {\small\textbf{בסיס הנתונים:} #2 \hfill \textbf{ניקוד:} #3}\par\vspace{2pt}%
  {\color{rulegray}\hrule height 0.4pt}\vspace{6pt}}

\newcommand{\rulebar}{\vspace{-4pt}{\color{rulegray}\hrule height 0.8pt}\vspace{5pt}}

% ---- שאלה וסעיף ----
\newcounter{qnum}
\newcommand{\question}[2]{%
  \needspace{5\baselineskip}\refstepcounter{qnum}\par\vspace{8pt}%
  \noindent\textbf{\color{inkblue}שאלה \theqnum.}\ \textbf{#2}\hfill{\small(#1 נק')}\par\vspace{2pt}}
\newcommand{\subq}[2]{\needspace{3\baselineskip}\par\vspace{4pt}\noindent\textbf{#1.}~#2\par}
\newcommand{\subqp}[3]{\needspace{8\baselineskip}\par\vspace{5pt}\noindent%
  \textbf{#1.}~{\small\textbf{\color{inkblue}(#3 נק')}}~#2\par}
\newcommand{\pairitem}[1]{\needspace{2\baselineskip}\par\vspace{4pt}\noindent\hspace{1.4em}\textbf{#1}\par}

% ---- שורות תשובה מנוקדות ----
\newcommand{\answerlines}[1]{%
  \par\vspace{1pt}%
  \needspace{\dimexpr 6mm*#1 + 2mm\relax}%
  \begingroup\color{rulegray}%
  \count255=#1\relax
  \loop\ifnum\count255>0
    \noindent\babelsublr{\rule[-2pt]{0.6pt}{9pt}\rule[-2pt]{\dimexpr\linewidth-0.6pt}{0.3pt}}\par\vspace{5.4mm}%
    \advance\count255 by -1
  \repeat
  \endgroup\vspace{-2mm}}

% ---- שורת תשובה קצרה אחת ----
\newcommand{\blankline}[1]{\babelsublr{\rule[-2pt]{#1}{0.4pt}}}

% ---- קוד SQL, תמיד משמאל לימין ----
\DefineVerbatimEnvironment{sqlverb}{Verbatim}{fontsize=\small, frame=leftline, framerule=1.2pt,
  rulecolor=\color{rulegray}, framesep=7pt, xleftmargin=12pt, samepage=true}
\DefineVerbatimEnvironment{sqlbad}{Verbatim}{fontsize=\small, frame=leftline, framerule=1.2pt,
  rulecolor=\color{warnred}, framesep=7pt, xleftmargin=12pt, samepage=true}
\DefineVerbatimEnvironment{sqlgood}{Verbatim}{fontsize=\small, frame=leftline, framerule=1.2pt,
  rulecolor=\color{okgreen}, framesep=7pt, xleftmargin=12pt, samepage=true}

\newcommand{\lr}[1]{\babelsublr{\mbox{#1}}}
\newcommand{\sqlin}[1]{\babelsublr{\texttt{#1}}}

% ---- תיבה אפורה ----
\newsavebox{\panelbox}
\newenvironment{panel}{\par\addvspace{7pt}\begin{lrbox}{\panelbox}%
  \begin{minipage}{\dimexpr\linewidth-2\fboxsep-2\fboxrule}}%
  {\end{minipage}\end{lrbox}\noindent\fcolorbox{rulegray}{boxgray}{\usebox{\panelbox}}\par\addvspace{7pt}}

% ---- תיבה צבעונית עם כותרת ----
\newsavebox{\notebox}
\newenvironment{notepanel}[2]{\def\npcolor{#1}\par\addvspace{7pt}\begin{lrbox}{\notebox}%
  \begin{minipage}{\dimexpr\linewidth-2\fboxsep-2\fboxrule}%
  \textbf{\color{#1}#2}\par\vspace{2pt}}%
  {\end{minipage}\end{lrbox}\noindent\fcolorbox{\npcolor}{white}{\usebox{\notebox}}\par\addvspace{7pt}}

% ---- טבלת נתונים מלאה ----
\newcommand{\thd}[1]{\textbf{#1}}
\newenvironment{fulltable}[3]{%
  \needspace{5\baselineskip}\par\addvspace{7pt}%
  \noindent{\small הטבלה \babelsublr{\texttt{\textbf{#1}}} \textendash\ #2#3}\par\vspace{2pt}%
  \begin{otherlanguage}{english}\begin{adjustbox}{max width=\linewidth}}
  {\end{adjustbox}\end{otherlanguage}\par\addvspace{4pt}}

% ---- פלט של שאילתה ----
\newenvironment{outtable}[1]{%
  \needspace{4\baselineskip}\par\vspace{2pt}%
  \noindent{\small\textbf{הפלט}#1}\par\vspace{1pt}%
  \begin{otherlanguage}{english}\begin{adjustbox}{max width=\linewidth}}
  {\end{adjustbox}\end{otherlanguage}\par\addvspace{5pt}}

% ---- פלט ללא כותרת ----
\newenvironment{outplain}{%
  \par\vspace{2pt}\begin{otherlanguage}{english}\begin{adjustbox}{max width=\linewidth}}
  {\end{adjustbox}\end{otherlanguage}\par\addvspace{4pt}}

\newcommand{\nullv}{\textit{\small NULL}}
\providecommand{\tightlist}{\setlength{\itemsep}{0pt}\setlength{\parskip}{0pt}}


\fancyhead[R]{\small בסיסי נתונים ו\lr{-}\lr{SQL} \textbar{} חומר לימוד לבחינת אמצע}
\fancyhead[L]{\small \lr{\texttt{SELECT} · \texttt{WHERE} · \texttt{ORDER BY}}}

\begin{document}

\begin{center}
{\LARGE\bfseries\color{inkblue} כל מה שצריך לדעת לבחינה}\\[4pt]
{\large טבלאות ומפתחות \textbullet{} \texttt{SELECT} \textbullet{} \texttt{WHERE} \textbullet{} \texttt{ORDER BY}}\\[6pt]
{\large\bfseries חומר לימוד ודף עזר}
\end{center}
\vspace{2pt}
{\color{inkblue}\hrule height 1.2pt}
\vspace{8pt}

\begin{panel}
\textbf{איך משתמשים בדף הזה}
\begin{itemize}[leftmargin=1.3em, itemsep=2pt, topsep=3pt]
  \item \textbf{זה כל החומר של הבחינה, ולא יותר.} כל מה שכתוב כאן \textendash{} נלמד בכיתה. כל מה שלא כתוב כאן (צירוף טבלאות, \lr{\texttt{COUNT}, \texttt{GROUP BY}}, שינוי נתונים) \textbf{לא ייבחן}.
  \item \textbf{מותר להביא את הדף הזה לבחינה}, מודפס או בכתב יד. אל תסתמכו עליו בלבד: מי שלא תרגל, לא ימצא בו את התשובה בזמן.
  \item \textbf{כל הדוגמאות רצו באמת.} הפלט שמופיע כאן הוא הפלט האמיתי של \lr{SQLite}, ולא כתוב ביד. אפשר להעתיק כל דוגמה ל-\lr{\texttt{onecompiler.com/sqlite}} ולראות את אותו דבר.
  \item \textbf{פרק \lr{8} הוא דף העזר} \textendash{} טבלה אחת עם כל התחביר. בזמן הבחינה תסתכלו בעיקר בו.
\end{itemize}
\end{panel}

\vspace{4pt}
\noindent\begin{minipage}{\linewidth}
{\small\textbf{מה יש בדף}}\par\vspace{3pt}
{\small
\begin{tabular}{@{}p{0.30\linewidth}p{0.30\linewidth}p{0.34\linewidth}@{}}
\toprule
\thd{היסודות} & \thd{השליפה} & \thd{לקראת הבחינה} \\
\midrule
\lr{1.} טבלאות, מפתחות וקשרים & \lr{3.\ \texttt{SELECT}}: מה מציגים & \lr{7.} מהשאלה בעברית לשאילתה \\
\lr{2.} מבנה השאילתה וסדר הביצוע & \lr{4.\ \texttt{NULL}}: הערך החסר & \lr{8.} דף העזר \textendash{} הכול בטבלה אחת \\
 & \lr{5.\ \texttt{WHERE}}: אילו שורות & \lr{9.} הטעויות שחוזרות על עצמן \\
 & \lr{6.\ \texttt{ORDER BY}}: באיזה סדר & \lr{10.} תרגול א': כתיבת שאילתות \\
 & & \lr{11.} תרגול ב': מסוג הבחינה \\
\bottomrule
\end{tabular}}
\end{minipage}

\vspace{10pt}
\section{טבלאות, מפתחות וקשרים}
\rulebar

בסיס נתונים הוא אוסף של \textbf{טבלאות}. בכל טבלה יש \textbf{עמודות} (שדות) \textendash{} אלה הן התכונות שאנחנו שומרים, ו\textbf{שורות} (רשומות) \textendash{} אלה הן המופעים עצמם.

\vspace{2pt}
\begin{panel}
\textbf{הדוגמה של הפרק הזה: ליגת כדורסל.} \lr{\texttt{Teams}} (קבוצות), \lr{\texttt{Players}} (שחקנים), \lr{\texttt{Towns}} (ערים), \lr{\texttt{Halls}} (אולמות) ו-\lr{\texttt{PlayerTour}} (מי משתתף באיזה טורניר). בטבלת \lr{\texttt{Players}} כל שורה היא שחקן אחד, וכל עמודה היא דבר אחד שאנחנו יודעים על כל שחקן.
\end{panel}

\subsection{מפתח ראשי \lr{(Primary Key)}}

\textbf{המפתח הראשי הוא העמודה שמזהה שורה אחת ויחידה.} שתי שורות לא יכולות להחזיק את אותו ערך בעמודה הזאת, והערך בה לעולם אינו \lr{\texttt{NULL}}.

\begin{panel}
\textbf{איך מזהים מפתח ראשי בבחינה?} מחפשים את העמודה שהערכים בה \textbf{לא חוזרים אף פעם}, ושכל העמודות האחרות "תלויות" בה. בדרך כלל זו העמודה הראשונה, ובדרך כלל יש בשמה \lr{\texttt{Num}}, \lr{\texttt{Code}} או \lr{\texttt{Id}}: \lr{\texttt{PlayerNum}}, \lr{\texttt{TownCode}}, \lr{\texttt{TeamCode}}.

\vspace{3pt}
\textbf{מפתח מורכב.} לפעמים אין עמודה אחת שמספיקה, ואז המפתח הראשי בנוי מ\textbf{שתי עמודות יחד}. זה המצב הרגיל בטבלה מקשרת: בטבלה שרושמת אילו שחקנים משתתפים באיזה טורניר, אותו שחקן מופיע בכמה שורות ואותו טורניר מופיע בכמה שורות \textendash{} אבל \textbf{הצירוף} של השניים מופיע פעם אחת בלבד, ולכן הוא המפתח.
\end{panel}

\subsection{מפתח זר \lr{(Foreign Key)}}

\textbf{מפתח זר הוא עמודה שמצביעה על המפתח הראשי של טבלה אחרת.} זה הדבק שמחבר את הטבלאות. בעמודה \lr{\texttt{TeamTown}} שבטבלת הקבוצות נשמר \textbf{קוד} של עיר, ולא שם של עיר \textendash{} והקוד הזה קיים בטבלת הערים, בעמודה \lr{\texttt{TownCode}}.

\begin{notepanel}{tipgold}{שימו לב \textendash{} זה בדיוק מה שמבלבל בבחינה}
לשני השדות יש \textbf{שמות שונים}: \lr{\texttt{TeamTown}} מצביע על \lr{\texttt{TownCode}}. בבחינה של תשפ"ז לכל שדה בבסיס הנתונים שם \textbf{ייחודי}, ולכן אי אפשר לזהות מפתח זר לפי שם זהה. מזהים אותו \textbf{לפי המשמעות ולפי הנתונים}: אם בעמודה יש מספרים קטנים שחוזרים על עצמם, והמספרים האלה הם בדיוק המפתח הראשי של טבלה אחרת \textendash{} זה מפתח זר.
\end{notepanel}

\vspace{2pt}
\textbf{שלוש שאלות שעוזרות להחליט:} (א) האם הערך בעמודה הוא "קוד של משהו"? (ב) האם יש טבלה שבה הקוד הזה הוא המפתח הראשי? (ג) האם הערכים חוזרים על עצמם? אם כן לשלושתן \textendash{} זה מפתח זר, והקשר הוא \lr{1:N}.

\subsection{שלושה סוגי קשרים}

\vspace{2pt}
\noindent\begin{tabular}{@{}p{0.26\linewidth}p{0.33\linewidth}p{0.34\linewidth}@{}}
\toprule
\thd{הקשר} & \thd{מה זה אומר} & \thd{איך מזהים בטבלאות} \\
\midrule
\textbf{\lr{1:1}} \newline אחד לאחד & לכל שורה בצד אחד מתאימה \textbf{שורה אחת לכל היותר} בצד השני & המפתח הזר הוא גם \textbf{ייחודי} \lr{(\texttt{UNIQUE})}. לדוגמה: לכל קבוצה אולם בית אחד \\
\textbf{\lr{1:N}} \newline אחד לרבים & לשורה אחת בצד ה"אחד" מתאימות \textbf{הרבה} שורות בצד ה"רבים" & המפתח הזר יושב בצד ה\textbf{רבים}, והערכים בו \textbf{חוזרים}. לדוגמה: קבוצה אחת, הרבה שחקנים \\
\textbf{\lr{M:N}} \newline רבים לרבים & לכל שורה בצד אחד הרבה שורות בצד השני, ו\textbf{להפך} & יש \textbf{טבלה שלישית}, מקשרת, שכל תוכנה הוא שני מפתחות זרים. לדוגמה: שחקן משתתף בכמה טורנירים, ובטורניר משתתפים הרבה שחקנים \\
\bottomrule
\end{tabular}

\vspace{6pt}
\begin{notepanel}{warnred}{הכלל שפותר את רוב השאלות}
\textbf{המפתח הזר תמיד יושב בצד ה"רבים".} בקשר "עיר אחת \textendash{} הרבה קבוצות", העמודה \lr{\texttt{TeamTown}} נמצאת בטבלת ה\textbf{קבוצות}, ולא בטבלת הערים. אם ראיתם עמודה שהערכים בה חוזרים ומצביעים לטבלה אחרת \textendash{} אתם בצד הרבים.
\end{notepanel}

\begin{notepanel}{tipgold}{שאלה שחוזרת: "יש עיר שאין בה אף אחד" \textendash{} האם זו תקלה?}
\textbf{לא. זה תקין לגמרי.} בקשר \lr{1:N}, לשורה בצד ה"אחד" מותר שלא יהיו לה \textbf{בכלל} שורות מתאימות בצד ה"רבים": עיר שרשומה בטבלת הערים, ואין בה אף קבוצה; קטגוריה שאין בה אף פריט. \textbf{ההפך אסור:} ערך במפתח הזר חייב להתאים לשורה קיימת \textendash{} אי אפשר לרשום קבוצה בעיר שאינה בטבלת הערים.

\vspace{3pt}
\textbf{איך בודקים בבחינה:} עוברים על העמודה של המפתח הזר, ומסמנים אילו קודים מופיעים בה. קוד שנמצא בטבלת העזר ולא מופיע אף פעם במפתח הזר \textendash{} הוא ה"ריק" שהשאלה מחפשת.
\end{notepanel}

\subsection{תרשים \lr{DSD}}

ב\lr{DSD} כל טבלה היא מלבן, שבתוכו שמות השדות. המפתח הראשי מסומן למעלה (או בקו תחתי), וכל מפתח זר מחובר ב\textbf{חץ} אל המפתח הראשי שהוא מצביע עליו. כך נראה תרשים של שלוש טבלאות:

\vspace{4pt}
\begin{otherlanguage}{english}
\begin{center}
\begin{tikzpicture}[font=\scriptsize\sffamily]
\tikzset{tblm/.style={matrix of nodes,
  nodes={anchor=west, inner sep=2.6pt, text width=25mm, align=left, font=\scriptsize\sffamily},
  row sep=-\pgflinewidth, draw=inkblue, line width=0.7pt, inner sep=0pt,
  row 1/.style={nodes={fill=inkblue, text=white, font=\scriptsize\sffamily\bfseries,
                       align=center, text width=25mm}}}}
\matrix (C) [tblm] { Towns \\ \underline{TownCode} \\ TownName \\ };
\matrix (T) [tblm, right=17mm of C] { Teams \\ \underline{TeamCode} \\ TeamName \\ TeamTown \\ FoundYear \\ };
\matrix (R) [tblm, right=17mm of T] { Players \\ \underline{PlayerNum} \\ PlayerName \\ PlayerTeam \\ BirthYear \\ };
\draw[-latex, inkblue, line width=0.8pt] (T-4-1.west)
      -- node[midway, above, font=\tiny\sffamily, text=inkblue] {1:N} (C-2-1.east);
\draw[-latex, inkblue, line width=0.8pt] (R-4-1.west)
      -- node[midway, above, font=\tiny\sffamily, text=inkblue] {1:N} (T-2-1.east);
\end{tikzpicture}
\end{center}
\end{otherlanguage}

\vspace{2pt}
\textbf{איך מציירים את החצים בבחינה.} עוברים שדה-שדה ושואלים "האם זה קוד של משהו אחר?". \lr{\texttt{TeamTown}} הוא קוד של עיר, ולכן חץ ממנו אל \lr{\texttt{TownCode}}. \lr{\texttt{PlayerTeam}} הוא קוד של קבוצה, ולכן חץ ממנו אל \lr{\texttt{TeamCode}}. \lr{\texttt{PlayerName}} הוא טקסט חופשי \textendash{} ואינו מחובר לשום דבר.

\newpage
\section{מבנה השאילתה וסדר הביצוע}
\rulebar

כל שאילתה שתכתבו בבחינה בנויה מארבעה חלקים, \textbf{תמיד בסדר הזה}. שניים חייבים להופיע, ושניים לא:

\begin{otherlanguage}{english}
\begin{sqlverb}
SELECT     BookTitle, BookPrice     -- חייב:  אילו עמודות להציג
FROM       Books                    -- חייב:  מאיזו טבלה
WHERE      BookPrice > 70           -- רשות:  אילו שורות להחזיר
ORDER BY   BookPrice DESC;          -- רשות:  באיזה סדר להציג
\end{sqlverb}
\end{otherlanguage}

\begin{notepanel}{warnred}{הסדר הזה אינו המלצה אלא חוק}
\lr{\texttt{WHERE}} שנכתב אחרי \lr{\texttt{ORDER BY}} הוא \textbf{שגיאת תחביר}, והשאילתה לא תרוץ בכלל. אם אתם מתבלבלים, זכרו את ארבע המילים בסדר הזה: \textbf{מה \textendash{} מאיפה \textendash{} מי \textendash{} באיזה סדר}.
\end{notepanel}

\subsection{סדר הביצוע הלוגי \textendash{} ולמה חשוב להכיר אותו}

מנוע בסיס הנתונים \textbf{אינו} מבצע את השאילתה בסדר שבו היא כתובה. הוא מבצע אותה כך:

\vspace{3pt}
\begin{otherlanguage}{english}
\begin{center}
\begin{tikzpicture}[font=\small\sffamily, node distance=7mm]
\tikzset{st/.style={draw=rulegray, line width=0.6pt, rounded corners=2pt, fill=boxgray,
         inner sep=4pt, minimum height=9mm, text width=25mm, align=center, font=\small\sffamily}}
\node[st] (d) {4.~\textbf{ORDER BY}};
\node[st, right=of d] (c) {3.~\textbf{SELECT}};
\node[st, right=of c] (b) {2.~\textbf{WHERE}};
\node[st, right=of b] (a) {1.~\textbf{FROM}};
\draw[-latex, line width=0.9pt, inkblue] (a) -- (b);
\draw[-latex, line width=0.9pt, inkblue] (b) -- (c);
\draw[-latex, line width=0.9pt, inkblue] (c) -- (d);
\end{tikzpicture}
\end{center}
\end{otherlanguage}
\vspace{-6pt}
\begin{center}
{\small קוראים את הטבלה \ \textbullet\ \ מסננים שורות \ \textbullet\ \ בוחרים עמודות \ \textbullet\ \ ממיינים את התוצאה}
\end{center}

\vspace{4pt}
\textbf{שתי מסקנות מעשיות שנובעות מכאן:}

\begin{enumerate}[leftmargin=1.5em, itemsep=3pt, topsep=3pt]
  \item \lr{\texttt{ORDER BY}} מתבצע \textbf{אחרי} \lr{\texttt{SELECT}}, ולכן \textbf{מותר} למיין לפי כינוי \lr{(AS)} שהגדרתם \textendash{} הכינוי כבר קיים.
  \item \lr{\texttt{WHERE}} מתבצע \textbf{לפני} \lr{\texttt{SELECT}}, ולכן \textbf{אסור} לסנן לפי כינוי \textendash{} בשלב הסינון הוא עוד לא נולד. בתנאי כותבים את החישוב במלואו.
\end{enumerate}

\begin{panel}
\textbf{וגם:} התוצאה של כל שאילתה היא \textbf{טבלה חדשה}, והטבלה המקורית \textbf{אינה משתנה}. \lr{\texttt{SELECT}} רק מציג \textendash{} הוא לא מוחק, לא ממיין ולא משנה שום דבר בבסיס הנתונים. גם \lr{\texttt{ORDER BY}} אינו ממיין את הטבלה: הוא ממיין רק את \textbf{הפלט} של השאילתה הזאת.
\end{panel}

\vspace{4pt}
\section{\lr{\texttt{SELECT}} \textendash{} מה מציגים}
\rulebar

כל הדוגמאות בפרקים \lr{3\textendash7} רצות על בסיס הנתונים הקטן הזה, \textbf{מועדון הקריאה}. שתי טבלאות, וקשר \lr{1:N} ביניהן. שימו לב ל\lr{\texttt{NULL}}-ים: לספר \lr{3} אין מדף, ולשני קוראים אין טלפון.

\begin{fulltable}{Books}{7 שורות}{}
\begin{tabular}{lllll}
\toprule
\thd{BookNum} & \thd{BookTitle} & \thd{BookPrice} & \thd{BookYear} & \thd{BookShelf} \\
\midrule
1 & מסע אל הירח & 68 & 2018 & A \\
2 & סודות הגלקסיה & 95 & 2021 & B \\
3 & יומן מהמחסן & 45 & 2018 & \nullv \\
4 & הדרך לירכא & 72 & 2023 & A \\
5 & ספר המתכונים & 120 & 2021 & C \\
6 & שירים מהגליל & 45 & 2019 & B \\
7 & אגדות הגליל & 38 & 2018 & A \\
\bottomrule
\end{tabular}
\end{fulltable}

\begin{fulltable}{Readers}{4 שורות}{}
\begin{tabular}{lllll}
\toprule
\thd{ReaderNum} & \thd{ReaderFirst} & \thd{ReaderLast} & \thd{JoinDate} & \thd{ReaderPhone} \\
\midrule
1 & נור & סאלם & 2024-09-01 & 050-1111111 \\
2 & איתי & ברק & 2025-02-10 & \nullv \\
3 & לינא & חדאד & 2025-09-01 & 052-3333333 \\
4 & עומר & לוי & 2023-11-15 & \nullv \\
\bottomrule
\end{tabular}
\end{fulltable}

\subsection{בחירת עמודות}

\begin{otherlanguage}{english}
\begin{sqlverb}
SELECT * FROM Books;                      -- כל העמודות, בסדר שבו הן בטבלה
SELECT BookTitle, BookPrice FROM Books;   -- רק שתי עמודות, בסדר שביקשתי
SELECT BookPrice, BookTitle FROM Books;   -- אותן עמודות, בסדר הפוך
\end{sqlverb}
\end{otherlanguage}

\textbf{הכוכבית נוחה אך מסוכנת.} בבחינה, שאלה שמבקשת "הצג את שם הספר והמחיר" דורשת שתי עמודות בדיוק. \lr{\texttt{SELECT *}} יציג חמש, וזו תשובה שגויה.

\subsection{כינויים \lr{(\texttt{AS})} \textendash{} לשנות את כותרת העמודה}

\begin{otherlanguage}{english}
\begin{sqlverb}
SELECT BookTitle AS "שם הספר", BookPrice AS "מחיר בשח"
FROM   Books;
\end{sqlverb}
\end{otherlanguage}

\begin{outtable}{ \textendash{} הכותרות השתנו, הנתונים לא}
\begin{tabular}{ll}
\toprule
\thd{שם הספר} & \thd{מחיר בשח} \\
\midrule
מסע אל הירח & 68 \\
סודות הגלקסיה & 95 \\
יומן מהמחסן & 45 \\
\multicolumn{2}{l}{\textit{\small ... ועוד 4 שורות}} \\
\bottomrule
\end{tabular}
\end{outtable}

\begin{notepanel}{warnred}{גרש בודד לערך, גרש כפול לשם}
זה הכלל שמפיל הכי הרבה תלמידים:
\begin{itemize}[leftmargin=1.3em, itemsep=1pt, topsep=2pt]
  \item \sqlin{'A'} \textendash{} עם גרש \textbf{בודד} \textendash{} הוא \textbf{ערך}: המחרוזת \lr{A}. כך כותבים טקסט בתוך \lr{\texttt{WHERE}}.
  \item \sqlin{"שם הספר"} \textendash{} עם גרש \textbf{כפול} \textendash{} הוא \textbf{שם}: של עמודה או של כינוי. כך כותבים כינוי שיש בו רווח או עברית.
\end{itemize}
\vspace{2pt}
לכן \sqlin{WHERE BookShelf = "A"} מחפש \textbf{עמודה} בשם \lr{A}. באורקל, ובבחינת הבגרות, זו \textbf{שגיאה}. ב-\lr{SQLite} יש "רשת ביטחון": אם אין עמודה בשם הזה, הוא מתייחס לזה כטקסט \textendash{} והשאילתה עובדת במקרה. \textbf{אל תסמכו על זה}: אותה שאילתה תיפול בבחינה. וכינוי של מילה אחת באנגלית, בלי רווח, לא דורש גרשיים בכלל: \sqlin{AS AfterSale}.
\end{notepanel}

\subsection{חישובים בתוך \lr{\texttt{SELECT}}}

מותר לחשב על עמודה, והחישוב מתבצע \textbf{לכל שורה בנפרד}. האופרטורים: \lr{\texttt{+}\ \ \texttt{-}\ \ \texttt{*}\ \ \texttt{/}}.

\begin{otherlanguage}{english}
\begin{sqlverb}
SELECT BookTitle, BookPrice, BookPrice * 0.9 AS AfterSale
FROM   Books;
\end{sqlverb}
\end{otherlanguage}

\begin{outplain}
\begin{tabular}{lll}
\toprule
\thd{BookTitle} & \thd{BookPrice} & \thd{AfterSale} \\
\midrule
מסע אל הירח & 68 & 61.2 \\
סודות הגלקסיה & 95 & 85.5 \\
יומן מהמחסן & 45 & 40.5 \\
הדרך לירכא & 72 & 64.8 \\
ספר המתכונים & 120 & 108.0 \\
שירים מהגליל & 45 & 40.5 \\
אגדות הגליל & 38 & 34.2 \\
\bottomrule
\end{tabular}
\end{outplain}

{\small\textbf{הוספת מע"מ \lr{18\%}:} מכפילים ב-\sqlin{1.18}. \textbf{הנחה של \lr{10\%}:} מכפילים ב-\sqlin{0.9}. \textbf{לשנה שלמה:} מכפילים ב-\sqlin{12}.}

\subsection{שרשור מחרוזות \lr{(\texttt{||})}}

שני קווים אנכיים \textbf{מדביקים} שתי מחרוזות. כדי לקבל רווח ביניהן, מדביקים גם מחרוזת של רווח:

\begin{otherlanguage}{english}
\begin{sqlverb}
SELECT ReaderFirst || ' ' || ReaderLast AS FullName
FROM   Readers;
\end{sqlverb}
\end{otherlanguage}

\begin{outplain}
\begin{tabular}{l}
\toprule
\thd{FullName} \\
\midrule
נור סאלם \\
איתי ברק \\
לינא חדאד \\
עומר לוי \\
\bottomrule
\end{tabular}
\end{outplain}

\subsection{\lr{\texttt{DISTINCT}} \textendash{} בלי כפילויות}

\lr{\texttt{DISTINCT}} נכתב \textbf{פעם אחת}, מיד אחרי \lr{\texttt{SELECT}}, ומסלק שורות פלט זהות.

\begin{otherlanguage}{english}
\begin{sqlverb}
SELECT DISTINCT BookYear FROM Books;     -- אילו שנים בכלל קיימות?
\end{sqlverb}
\end{otherlanguage}

\begin{outplain}
\begin{tabular}{l}
\toprule
\thd{BookYear} \\
\midrule
2018 \\
2021 \\
2023 \\
2019 \\
\bottomrule
\end{tabular}
\end{outplain}

{\small מ-\lr{7} שורות נשארו \lr{4}: \lr{2018} הופיעה שלוש פעמים ו-\lr{2021} פעמיים. שימו לב שהתוצאה \textbf{אינה ממוינת} \textendash{} \lr{\texttt{DISTINCT}} מסלק כפילויות, ולא ממיין. רוצים גם סדר? הוסיפו \lr{\texttt{ORDER BY}}.}

\begin{notepanel}{tipgold}{\lr{\texttt{DISTINCT}} עובד על \textbf{השורה} כולה, לא על עמודה אחת}
\sqlin{SELECT DISTINCT BookYear, BookShelf FROM Books} יחזיר \lr{6} שורות ולא \lr{4}: הצירוף \lr{(2018, A)} מופיע פעמיים ולכן נספר פעם אחת, אבל \lr{(2021, B)} ו-\lr{(2021, C)} הם \textbf{צירופים שונים} ושניהם נשארים. השאלה היא תמיד "האם \textbf{כל השורה} בפלט חוזרת?".
\end{notepanel}

\newpage
\section{\lr{\texttt{NULL}} \textendash{} הערך החסר}
\rulebar

\lr{\texttt{NULL}} הוא \textbf{לא אפס}, \textbf{לא מחרוזת ריקה} ו\textbf{לא רווח}. הוא אומר דבר אחד: \textbf{אנחנו לא יודעים מה הערך}. לספר \lr{3} אין מדף רשום, ולאיתי ולעומר אין טלפון.

\subsection{\lr{\texttt{NULL}} "מדביק"}

כל חישוב שבו משתתף \lr{\texttt{NULL}} מחזיר \lr{\texttt{NULL}} \textendash{} גם חשבון, וגם שרשור מחרוזות:

\begin{otherlanguage}{english}
\begin{sqlverb}
SELECT ReaderFirst || ' ' || ReaderPhone AS Res FROM Readers;
\end{sqlverb}
\end{otherlanguage}

\begin{outplain}
\begin{tabular}{l}
\toprule
\thd{Res} \\
\midrule
נור 050-1111111 \\
\nullv \\
לינא 052-3333333 \\
\nullv \\
\bottomrule
\end{tabular}
\end{outplain}

{\small השם \textbf{נעלם} אצל איתי ועומר. לא כתוב "איתי" ואחריו כלום \textendash{} \textbf{כל הביטוי} הפך ל-\lr{\texttt{NULL}}. זו שאלת "השלמת ערכים" קלאסית בבחינה.}

\subsection{איך בודקים \lr{\texttt{NULL}}: \lr{\texttt{IS NULL}}}

\begin{notepanel}{warnred}{\lr{\texttt{= NULL}} לא עובד. אף פעם.}
\sqlin{WHERE BookShelf = NULL} מחזיר \textbf{אפס שורות}, גם כשיש שורות בלי מדף \textendash{} והוא \textbf{אינו} שגיאת תחביר, ולכן הוא לא מתריע. הסיבה: אי אפשר להשוות לערך שלא ידוע. התשובה ל"האם הלא-ידוע שווה ל-\lr{5}?" היא "לא ידוע", ו"לא ידוע" אינו "אמת" \textendash{} ולכן השורה לא נכנסת לפלט.

\vspace{3pt}
\textbf{הדרך היחידה:} \sqlin{IS NULL} \ ו-\ \sqlin{IS NOT NULL}.
\end{notepanel}

\begin{otherlanguage}{english}
\begin{sqlverb}
SELECT BookNum, BookTitle, BookShelf FROM Books WHERE BookShelf IS NULL;
SELECT ReaderFirst, ReaderPhone     FROM Readers WHERE ReaderPhone IS NOT NULL;
\end{sqlverb}
\end{otherlanguage}

\vspace{2pt}
\noindent\begin{minipage}[t]{0.48\linewidth}
\begin{outplain}
\begin{tabular}{lll}
\toprule
\thd{BookNum} & \thd{BookTitle} & \thd{BookShelf} \\
\midrule
3 & יומן מהמחסן & \nullv \\
\bottomrule
\end{tabular}
\end{outplain}
\end{minipage}\hfill
\begin{minipage}[t]{0.48\linewidth}
\begin{outplain}
\begin{tabular}{ll}
\toprule
\thd{ReaderFirst} & \thd{ReaderPhone} \\
\midrule
נור & 050-1111111 \\
לינא & 052-3333333 \\
\bottomrule
\end{tabular}
\end{outplain}
\end{minipage}

\vspace{4pt}
\begin{panel}
\textbf{ההשלכה הנסתרת, ששווה נקודות בבחינה.} שורה עם \lr{\texttt{NULL}} \textbf{אינה נכנסת גם לתנאי ההפוך}. התנאי \sqlin{WHERE BookShelf <> 'A'} \textbf{לא} יחזיר את ספר \lr{3}, למרות שהמדף שלו בוודאות אינו \lr{A} \textendash{} כי הוא לא ידוע. רוצים לכלול אותו? כותבים במפורש:

\vspace{2pt}
\sqlin{WHERE BookShelf <> 'A' OR BookShelf IS NULL}
\end{panel}

\newpage
\section{\lr{\texttt{WHERE}} \textendash{} אילו שורות}
\rulebar

\lr{\texttt{WHERE}} הוא \textbf{מסננת}. המנוע עובר שורה-שורה, בודק את התנאי, ומכניס לפלט רק את השורות שבהן התנאי יצא \textbf{אמת}.

\subsection{אופרטורי השוואה}

\vspace{2pt}
\noindent\begin{tabular}{@{}p{0.16\linewidth}p{0.30\linewidth}p{0.46\linewidth}@{}}
\toprule
\thd{האופרטור} & \thd{המשמעות} & \thd{דוגמה} \\
\midrule
\sqlin{=} & שווה & \sqlin{WHERE BookYear = 2018} \\
\sqlin{<>} & \textbf{אינו} שווה & \sqlin{WHERE BookShelf <> 'A'} \\
\lr{\texttt{>} \ \texttt{<}} & גדול / קטן & \sqlin{WHERE BookPrice > 70} \\
\lr{\texttt{>=} \ \texttt{<=}} & גדול או שווה / קטן או שווה & \sqlin{WHERE BookPrice >= 45} \\
\bottomrule
\end{tabular}

\vspace{5pt}
{\small\textbf{שימו לב:} "שווה" הוא סימן \textbf{אחד} (\sqlin{=}), ולא שניים. "אינו שווה" הוא \sqlin{<>} \textendash{} קטן ואחריו גדול. (\sqlin{!=} עובד ברוב המנועים, אבל בבחינה כתבו \sqlin{<>}.)}

\subsection{טקסט ותאריכים}

\begin{otherlanguage}{english}
\begin{sqlverb}
WHERE BookShelf  = 'A'                 -- טקסט: גרש בודד
WHERE JoinDate  >= '2025-01-01'        -- תאריך: גם הוא טקסט, בגרש בודד
\end{sqlverb}
\end{otherlanguage}

\textbf{למה השוואת תאריכים בכלל עובדת?} כי התאריך נשמר בפורמט \lr{\texttt{YYYY-MM-DD}} \textendash{} שנה, אחריה חודש, אחריה יום. בפורמט הזה הסדר ה\textbf{אלפביתי} של המחרוזות הוא בדיוק הסדר ה\textbf{כרונולוגי}, ולכן \sqlin{'2026-03-02' < '2026-05-15'} הוא אמת. זו הסיבה שלא כותבים תאריך כ-\lr{\texttt{02/03/2026}}.

\subsection{\lr{\texttt{BETWEEN}} \textendash{} טווח}

\begin{otherlanguage}{english}
\begin{sqlverb}
SELECT BookTitle, BookPrice FROM Books WHERE BookPrice BETWEEN 45 AND 72;
\end{sqlverb}
\end{otherlanguage}

\begin{outplain}
\begin{tabular}{ll}
\toprule
\thd{BookTitle} & \thd{BookPrice} \\
\midrule
מסע אל הירח & 68 \\
יומן מהמחסן & 45 \\
הדרך לירכא & 72 \\
שירים מהגליל & 45 \\
\bottomrule
\end{tabular}
\end{outplain}

\begin{notepanel}{warnred}{\lr{\texttt{BETWEEN}} \textbf{כולל} את שני הקצוות}
בפלט למעלה מופיעים גם \lr{45} וגם \lr{72}. \sqlin{BETWEEN 45 AND 72} הוא בדיוק \sqlin{BookPrice >= 45 AND BookPrice <= 72}.

\vspace{3pt}
\textbf{וגם:} הקטן תמיד נכתב \textbf{קודם}. \sqlin{BETWEEN 72 AND 45} יחזיר \textbf{אפס שורות} \textendash{} בלי שגיאה, ובלי אזהרה.
\end{notepanel}

\textbf{הכי שימושי לתאריכים}, כדי לקבל חודש או טווח חודשים שלם:

\begin{otherlanguage}{english}
\begin{sqlverb}
WHERE InDate BETWEEN '2026-03-01' AND '2026-05-31'    -- מרץ עד מאי, כולל
\end{sqlverb}
\end{otherlanguage}

\subsection{\lr{\texttt{IN}} \textendash{} רשימה של ערכים}

\lr{\texttt{IN}} מקצר שרשרת של \lr{\texttt{OR}}. שתי השורות האלה עושות אותו דבר בדיוק:

\begin{otherlanguage}{english}
\begin{sqlverb}
WHERE BookYear = 2019 OR BookYear = 2023      -- ארוך
WHERE BookYear IN (2019, 2023)                -- קצר, ועדיף
\end{sqlverb}
\end{otherlanguage}

\begin{outplain}
\begin{tabular}{ll}
\toprule
\thd{BookTitle} & \thd{BookYear} \\
\midrule
הדרך לירכא & 2023 \\
שירים מהגליל & 2019 \\
\bottomrule
\end{tabular}
\end{outplain}

{\small ברשימה של טקסט, כל ערך בגרש בודד: \sqlin{WHERE Priority IN ('דחוף', 'מיידי')}. ההפך: \sqlin{NOT IN (...)}.}

\subsection{\lr{\texttt{LIKE}} \textendash{} חיפוש לפי תבנית}

\lr{\texttt{LIKE}} מחפש \textbf{צורה} של מחרוזת, ולא התאמה מדויקת. שני סימנים מיוחדים:

\vspace{2pt}
\noindent\begin{tabular}{@{}p{0.13\linewidth}p{0.36\linewidth}p{0.43\linewidth}@{}}
\toprule
\thd{הסימן} & \thd{מה הוא מחליף} & \thd{דוגמה} \\
\midrule
\sqlin{\%} & \textbf{כל מספר} של תווים \textendash{} גם אפס & \sqlin{'ס\%'} = מתחיל ב-ס \\
\sqlin{\_} & \textbf{תו אחד} בדיוק & \sqlin{'\_\_\_'} = בדיוק שלושה תווים \\
\bottomrule
\end{tabular}

\vspace{4pt}
\noindent\begin{tabular}{@{}p{0.30\linewidth}p{0.62\linewidth}@{}}
\toprule
\thd{התבנית} & \thd{מה היא מוצאת} \\
\midrule
\sqlin{LIKE 'גל\%'} & \textbf{מתחיל} ב"גל" \\
\sqlin{LIKE '\%גל'} & \textbf{מסתיים} ב"גל" \\
\sqlin{LIKE '\%גל\%'} & \textbf{מכיל} "גל" \textendash{} בכל מקום במחרוזת \\
\bottomrule
\end{tabular}

\begin{otherlanguage}{english}
\begin{sqlverb}
SELECT BookNum, BookTitle FROM Books WHERE BookTitle LIKE '%גל%';
\end{sqlverb}
\end{otherlanguage}

\begin{outplain}
\begin{tabular}{ll}
\toprule
\thd{BookNum} & \thd{BookTitle} \\
\midrule
2 & סודות הגלקסיה \\
6 & שירים מהגליל \\
7 & אגדות הגליל \\
\bottomrule
\end{tabular}
\end{outplain}

\begin{notepanel}{tipgold}{שתי מלכודות ב-\lr{\texttt{LIKE}}}
\lr{1.} \textbf{בלי אחוז, \lr{\texttt{LIKE}} הוא פשוט \lr{\texttt{=}}.} \sqlin{LIKE 'גל'} ימצא רק מחרוזת שכולה "גל", וכאן \textendash{} אפס שורות.
\lr{2.} \textbf{"מכיל" דורש \textbf{שני} סימני אחוז}, אחד מכל צד. שאלה בנוסח "כל הספרים שמופיעה בכותרתם המילה גליל", כשהכותרת שמורה כולה בעמודה אחת, פותרים ב-\sqlin{LIKE '\%גליל\%'}.
\end{notepanel}

\subsection{\lr{\texttt{AND}}, \lr{\texttt{OR}}, \lr{\texttt{NOT}} \textendash{} ושאלת הסוגריים}

\begin{itemize}[leftmargin=1.3em, itemsep=1pt, topsep=2pt]
  \item \lr{\texttt{AND}} \textendash{} \textbf{שני} התנאים חייבים להתקיים. מצמצם את הפלט.
  \item \lr{\texttt{OR}} \textendash{} \textbf{מספיק אחד}. מרחיב את הפלט.
  \item \lr{\texttt{NOT}} \textendash{} הופך את התנאי.
\end{itemize}

\vspace{4pt}
\begin{notepanel}{warnred}{\lr{\texttt{AND}} גובר על \lr{\texttt{OR}} \textendash{} וזאת הטעות היקרה ביותר בבחינה}
בדיוק כמו שכפל גובר על חיבור בחשבון, \lr{\texttt{AND}} מתבצע \textbf{לפני} \lr{\texttt{OR}}. "ספרים ממדף \lr{A} או \lr{B}, שמחירם מעל \lr{60}" נראה כך \textendash{} ושימו לב להבדל בפלט:
\end{notepanel}

\vspace{-2pt}
\noindent\begin{minipage}[t]{0.49\linewidth}
\begin{otherlanguage}{english}
\begin{sqlbad}
WHERE BookShelf = 'A'
   OR BookShelf = 'B'
  AND BookPrice > 60
\end{sqlbad}
\end{otherlanguage}
\begin{outplain}
\begin{tabular}{lll}
\toprule
\thd{\#} & \thd{BookTitle} & \thd{BookPrice} \\
\midrule
1 & מסע אל הירח & 68 \\
2 & סודות הגלקסיה & 95 \\
4 & הדרך לירכא & 72 \\
\rowcolor{boxgray} 7 & אגדות הגליל & 38 \\
\bottomrule
\end{tabular}
\end{outplain}
{\scriptsize\textbf{\lr{4} שורות.} ספר \lr{7} נכנס למרות שמחירו \lr{38}: המנוע קרא "מדף \lr{A}" \textbf{או} "מדף \lr{B} ומעל \lr{60}".}
\end{minipage}\hfill
\begin{minipage}[t]{0.49\linewidth}
\begin{otherlanguage}{english}
\begin{sqlgood}
WHERE (BookShelf = 'A'
    OR BookShelf = 'B')
  AND BookPrice > 60
\end{sqlgood}
\end{otherlanguage}
\begin{outplain}
\begin{tabular}{lll}
\toprule
\thd{\#} & \thd{BookTitle} & \thd{BookPrice} \\
\midrule
1 & מסע אל הירח & 68 \\
2 & סודות הגלקסיה & 95 \\
4 & הדרך לירכא & 72 \\
\bottomrule
\end{tabular}
\end{outplain}
{\scriptsize\textbf{\lr{3} שורות} \textendash{} וזו התשובה הנכונה. הסוגריים כופים "קודם המדף, ואחר כך המחיר".}
\end{minipage}

\vspace{6pt}
\begin{panel}
\textbf{הכלל המעשי:} בכל פעם שיש \lr{\texttt{OR}} בשאילתה, \textbf{ועוד תנאי כלשהו} \textendash{} שימו סוגריים סביב ה-\lr{\texttt{OR}}. גם אם לא היו הכרחיים, הם לא יזיקו, והם מראים למעריך בדיוק מה התכוונתם. \textbf{עברית של שאלה שדורשת סוגריים:} "הספרים שנמצאים \textbf{במדף} \lr{A} \textbf{או במדף} \lr{B}, \textbf{ו}שמחירם מעל \lr{60}" \textendash{} ה"או" מתייחס רק למדף, ולכן הוא נכנס לסוגריים.
\end{panel}

\newpage
\section{\lr{\texttt{ORDER BY}} \textendash{} באיזה סדר}
\rulebar

\lr{\texttt{ORDER BY}} הוא \textbf{תמיד החלק האחרון} בשאילתה, והוא ממיין את הפלט. בלעדיו, הסדר שתקבלו אינו מובטח \textendash{} גם אם הוא נראה "הגיוני".

\subsection{עולה ויורד}

\begin{otherlanguage}{english}
\begin{sqlverb}
ORDER BY BookPrice           -- עולה: מהקטן לגדול   (ברירת המחדל)
ORDER BY BookPrice ASC       -- אותו דבר בדיוק, רק כתוב במפורש
ORDER BY BookPrice DESC      -- יורד: מהגדול לקטן
\end{sqlverb}
\end{otherlanguage}

{\small \lr{\texttt{ASC}} הוא ברירת המחדל, ולכן אין חובה לכתוב אותו. \lr{\texttt{DESC}} \textendash{} יש חובה. בטקסט, "עולה" הוא לפי סדר האלף-בית, ובתאריך \textendash{} מהמוקדם למאוחר.}

\begin{notepanel}{warnred}{\lr{\texttt{ORDER BY}} \textbf{אינו מסנן} \textendash{} מספר השורות לא משתנה}
זו שאלה אהובה בבחינה: "כמה שורות תחזיר השאילתה?". אם אין \lr{\texttt{WHERE}}, התשובה היא תמיד \textbf{כל השורות בטבלה} \textendash{} גם אם יש \lr{\texttt{ORDER BY}}, וגם אם יש \lr{\texttt{NULL}}-ים בעמודת המיון. \lr{\texttt{ORDER BY}} רק \textbf{מסדר מחדש} את אותן שורות.

\vspace{3pt}
בטבלת הספרים יש \lr{7} שורות. \sqlin{SELECT BookTitle FROM Books ORDER BY BookShelf} יחזיר \textbf{\lr{7}} שורות \textendash{} כולל הספר שאין לו מדף. \textbf{מי שמסנן \textendash{} זה \lr{\texttt{WHERE}}, והוא לבדו.}
\end{notepanel}

\subsection{מיון לפי כמה עמודות}

כותבים אותן \textbf{מופרדות בפסיק}. המנוע ממיין לפי הראשונה, ו\textbf{רק כששתי שורות זהות בה} \textendash{} הוא פונה לשנייה כדי להכריע. לכל עמודה כיוון משלה:

\begin{otherlanguage}{english}
\begin{sqlverb}
SELECT BookTitle, BookYear, BookPrice
FROM   Books
ORDER BY BookYear DESC, BookPrice;        -- שנה יורד, ובתוך שנה מחיר עולה
\end{sqlverb}
\end{otherlanguage}

\begin{outplain}
\begin{tabular}{lll}
\toprule
\thd{BookTitle} & \thd{BookYear} & \thd{BookPrice} \\
\midrule
הדרך לירכא & 2023 & 72 \\
סודות הגלקסיה & 2021 & 95 \\
ספר המתכונים & 2021 & 120 \\
שירים מהגליל & 2019 & 45 \\
אגדות הגליל & 2018 & 38 \\
יומן מהמחסן & 2018 & 45 \\
מסע אל הירח & 2018 & 68 \\
\bottomrule
\end{tabular}
\end{outplain}

{\small השנים יורדות: \lr{2023}, \lr{2021}, \lr{2019}, \lr{2018}. \textbf{בתוך} \lr{2021} המחירים עולים (\lr{95} ואז \lr{120}), ו\textbf{בתוך} \lr{2018} גם כן (\lr{38}, \lr{45}, \lr{68}). ה-\lr{\texttt{DESC}} חל \textbf{רק} על \lr{\texttt{BookYear}} \textendash{} כל עמודה צריכה \lr{\texttt{DESC}} משלה.}

\subsection{למיין לפי כינוי, ביטוי או מספר עמודה}

שלוש הדרכים האלה תקפות, ושלושתן מחזירות אותו פלט:

\begin{otherlanguage}{english}
\begin{sqlverb}
SELECT BookTitle, BookPrice * 0.9 AS AfterSale FROM Books ORDER BY AfterSale DESC;
SELECT BookTitle, BookPrice * 0.9 AS AfterSale FROM Books ORDER BY BookPrice * 0.9 DESC;
SELECT BookTitle, BookPrice * 0.9 AS AfterSale FROM Books ORDER BY 2 DESC;
\end{sqlverb}
\end{otherlanguage}

\begin{outplain}
\begin{tabular}{ll}
\toprule
\thd{BookTitle} & \thd{AfterSale} \\
\midrule
ספר המתכונים & 108.0 \\
סודות הגלקסיה & 85.5 \\
הדרך לירכא & 64.8 \\
מסע אל הירח & 61.2 \\
יומן מהמחסן & 40.5 \\
שירים מהגליל & 40.5 \\
אגדות הגליל & 34.2 \\
\bottomrule
\end{tabular}
\end{outplain}

{\small \lr{\texttt{ORDER BY 2}} פירושו "לפי העמודה ה\textbf{שנייה} ברשימת ה-\lr{\texttt{SELECT}}". זה עובד, אבל שביר: מי שיוסיף עמודה באמצע ישבור את המיון \textendash{} ולכן עדיף למיין לפי כינוי. \textbf{וזכרו את ההבדל מפרק \lr{2}:} ב-\lr{\texttt{ORDER BY}} \textbf{מותר} להשתמש בכינוי, וב-\lr{\texttt{WHERE}} \textbf{אסור}.}

\subsection{מיון טקסט \textendash{} שתי הפתעות}

\textbf{עברית ממוינת לפי האלף-בית}, כמו שמצפים: \sqlin{ORDER BY BookTitle} מחזיר אגדות, הדרך, יומן, מסע, סודות, ספר, שירים.

\vspace{3pt}
\begin{notepanel}{warnred}{אלפביתי \textendash{} ולא "הגיוני"}
המחשב אינו מבין את \textbf{המשמעות} של המילים, אלא רק משווה אותיות. לכן בעמודה של סוגי מנות, \textbf{"מנה עיקרית" תופיע לפני "מנה ראשונה"} \textendash{} שתיהן מתחילות ב"מנה~", וההכרעה נופלת על האות הבאה: \textbf{ע} לפני \textbf{ר}. זה נכון, גם אם בתפריט המנה הראשונה מוגשת קודם.

\vspace{3pt}
אותו דבר בדרגות, בגדלים ובימי השבוע: \textbf{אל תצפו לסדר ההגיוני} \textendash{} צפו לסדר האלף-בית.
\end{notepanel}

\begin{notepanel}{tipgold}{מספר שמאוחסן כטקסט ממוין כמו \textbf{מילה}}
אם עמודה מכילה טקסט, \lr{\texttt{'10'}} בא \textbf{לפני} \lr{\texttt{'9'}} \textendash{} כי ההשוואה היא תו-תו, והתו \lr{\texttt{'1'}} קטן מהתו \lr{\texttt{'9'}}. אותה השוואה על \textbf{מספרים} נותנת את התשובה ההפוכה:

\vspace{2pt}
\sqlin{SELECT '10' < '9'} \ \ מחזיר \ \textbf{\lr{1}} (אמת) \qquad \sqlin{SELECT 10 < 9} \ \ מחזיר \ \textbf{\lr{0}} (שקר)

\vspace{2pt}
לכן עמודה של מחירים או קודים צריכה להיות \lr{\texttt{INTEGER}}, ולא \lr{\texttt{TEXT}}.
\end{notepanel}

\subsection{\lr{\texttt{NULL}} במיון \textendash{} איפה הוא נופל}

\begin{otherlanguage}{english}
\begin{sqlverb}
SELECT BookNum, BookShelf FROM Books ORDER BY BookShelf;         -- עולה
SELECT BookNum, BookShelf FROM Books ORDER BY BookShelf DESC;    -- יורד
\end{sqlverb}
\end{otherlanguage}

\vspace{2pt}
\noindent\begin{minipage}[t]{0.47\linewidth}
\begin{outplain}
\begin{tabular}{ll}
\toprule
\thd{BookNum} & \thd{BookShelf} \\
\midrule
\rowcolor{boxgray} 3 & \nullv \\
1 & A \\
4 & A \\
7 & A \\
2 & B \\
6 & B \\
5 & C \\
\bottomrule
\end{tabular}
\end{outplain}
{\scriptsize עולה \textendash{} ה-\lr{\texttt{NULL}} \textbf{ראשון}.}
\end{minipage}\hfill
\begin{minipage}[t]{0.47\linewidth}
\begin{outplain}
\begin{tabular}{ll}
\toprule
\thd{BookNum} & \thd{BookShelf} \\
\midrule
5 & C \\
2 & B \\
6 & B \\
1 & A \\
4 & A \\
7 & A \\
\rowcolor{boxgray} 3 & \nullv \\
\bottomrule
\end{tabular}
\end{outplain}
{\scriptsize יורד \textendash{} ה-\lr{\texttt{NULL}} \textbf{אחרון}.}
\end{minipage}

\vspace{5pt}
\begin{panel}
\textbf{הכלל ב-\lr{SQLite}} (וזה המנוע שבו אנחנו עובדים): ב-\lr{\texttt{ASC}} ה-\lr{\texttt{NULL}}-ים ראשונים, וב-\lr{\texttt{DESC}} אחרונים \textendash{} כאילו \lr{\texttt{NULL}} הוא "הערך הקטן מכולם". \textbf{באורקל זה הפוך.} בבחינה, אם שאלה תלויה במקום של ה-\lr{\texttt{NULL}} \textendash{} סננו אותו במפורש ב-\lr{\texttt{WHERE}}, ואל תסתמכו על המיון.
\end{panel}

\subsection{\lr{\texttt{LIMIT}} \textendash{} "שלושת הראשונים"}

\lr{\texttt{LIMIT}} מחזיר רק את \lr{N} השורות הראשונות של הפלט, והוא נכתב \textbf{אחרי} \lr{\texttt{ORDER BY}}:

\begin{otherlanguage}{english}
\begin{sqlverb}
SELECT   BookTitle, BookPrice
FROM     Books
ORDER BY BookPrice DESC
LIMIT    3;                      -- שלושת הספרים היקרים ביותר
\end{sqlverb}
\end{otherlanguage}

\begin{outplain}
\begin{tabular}{ll}
\toprule
\thd{BookTitle} & \thd{BookPrice} \\
\midrule
ספר המתכונים & 120 \\
סודות הגלקסיה & 95 \\
הדרך לירכא & 72 \\
\bottomrule
\end{tabular}
\end{outplain}

\begin{notepanel}{warnred}{\lr{\texttt{LIMIT}} בלי \lr{\texttt{ORDER BY}} הוא חסר משמעות}
"שלושת היקרים" = \textbf{קודם למיין} לפי מחיר יורד, \textbf{ואחר כך} לקצוץ ל-\lr{3}. \sqlin{LIMIT 3} לבדו יחזיר שלוש שורות שרירותיות. הכיוון חייב להתאים לשאלה: היקרים \textendash{} \lr{\texttt{DESC}}; הזולים \textendash{} \lr{\texttt{ASC}}.
\end{notepanel}

\subsection{הכול ביחד}

\begin{otherlanguage}{english}
\begin{sqlverb}
SELECT   BookTitle, BookPrice          -- 3. מה מציגים
FROM     Books                         -- 1. מאיפה
WHERE    BookYear <> 2018              -- 2. אילו שורות
ORDER BY BookPrice DESC                -- 4. באיזה סדר
LIMIT    2;                            -- 5. כמה
\end{sqlverb}
\end{otherlanguage}

\begin{outplain}
\begin{tabular}{ll}
\toprule
\thd{BookTitle} & \thd{BookPrice} \\
\midrule
ספר המתכונים & 120 \\
סודות הגלקסיה & 95 \\
\bottomrule
\end{tabular}
\end{outplain}

\newpage
\section{מהשאלה בעברית לשאילתה}
\rulebar

רוב התלמידים שיודעים את התחביר עדיין מאבדים נקודות, מסיבה אחת: הם לא מצליחים לתרגם את \textbf{המשפט בעברית} לשאילתה. הפרק הזה הוא השיטה.

\subsection{ארבעה צעדים, תמיד באותו סדר}

\begin{panel}
\begin{enumerate}[leftmargin=1.6em, itemsep=3pt, topsep=2pt]
  \item \textbf{מאיזו טבלה?} חפשו במשפט את שם הדבר שעליו שואלים \textendash{} ספרים? קוראים? \ זה ה-\lr{\texttt{FROM}}.
  \item \textbf{אילו עמודות להציג?} סמנו במשפט את המילים שאחרי "\textbf{הציגו את}". \textbf{בדיוק אותן}, ולא יותר \textendash{} זה ה-\lr{\texttt{SELECT}}.
  \item \textbf{אילו שורות?} חפשו את המילה "\textbf{ש}" \textendash{} "ספרים \textbf{ש}מחירם...", "קוראים \textbf{ש}אין להם..." \ כל מה שבא אחריה הוא ה-\lr{\texttt{WHERE}}.
  \item \textbf{באיזה סדר?} חפשו "ממוינים", "מה... ל...", "הראשונים" \textendash{} זה ה-\lr{\texttt{ORDER BY}}. אם אין \textendash{} פשוט אין.
\end{enumerate}
\end{panel}

\subsection{המילון: מה כתוב בשאלה, ומה כותבים}

{\small שלוש טבלאות קצרות, לפי שלושת החלקים של השאילתה. זה הפרק שכדאי לסמן ולחזור אליו.}

\vspace{5pt}
\noindent\textbf{\color{inkblue}א. העמודות \textendash{} מה מציגים}\par\vspace{2pt}
\noindent\begin{tabular}{@{}p{0.46\linewidth}p{0.48\linewidth}@{}}
\toprule
\thd{מה כתוב בשאלה} & \thd{מה כותבים} \\
\midrule
"הציגו את שם הספר ואת המחיר" & \sqlin{SELECT BookTitle, BookPrice} \\
"תחת הכותרת \lr{X}" \ \textbar\ \ "בכותרת \lr{X}" & \sqlin{AS X} \ (עברית או רווח: גרש כפול) \\
"כל סוג \textbf{פעם אחת}" \ \textbar\ \ "אילו סוגים קיימים" & \sqlin{SELECT DISTINCT} \\
"המחיר כולל מע"מ \lr{18\%}" & \sqlin{BookPrice * 1.18} \\
"אחרי הנחה של \lr{10\%}" & \sqlin{BookPrice * 0.9} \\
"שם מלא" (שם פרטי ומשפחה יחד) & \sqlin{First || ' ' || Last} \\
\bottomrule
\end{tabular}

\vspace{7pt}
\noindent\textbf{\color{inkblue}ב. התנאי \textendash{} אילו שורות}\par\vspace{2pt}
\noindent\begin{tabular}{@{}p{0.46\linewidth}p{0.48\linewidth}@{}}
\toprule
\thd{מה כתוב בשאלה} & \thd{מה כותבים} \\
\midrule
"\textbf{מעל} \lr{30}" \ \textbar\ \ "גבוה מ-\lr{30}" & \sqlin{> 30} \ \ \textendash{} \lr{30} עצמו \textbf{לא} נכנס \\
"\lr{30} \textbf{ומעלה}" \ \textbar\ \ "לפחות \lr{30}" & \sqlin{>= 30} \ \ \textendash{} \lr{30} עצמו \textbf{כן} נכנס \\
"מתחת ל-\lr{30}" \ \textbar\ \ "זול מ-\lr{30}" & \sqlin{< 30} \\
"בין \lr{10} ל-\lr{20}, כולל" & \sqlin{BETWEEN 10 AND 20} \\
"אחד מ: א, ב, ג" & \sqlin{IN ('א', 'ב', 'ג')} \\
"שאינו \lr{X}" \ \textbar\ \ "חוץ מ-\lr{X}" & \sqlin{<> 'X'} \\
"\textbf{מכיל}" \ \textbar\ \ "מופיעה בו המילה" & \sqlin{LIKE '\%...\%'} \ \ \textendash{} \textbf{שני} אחוזים \\
"מתחיל ב" & \sqlin{LIKE '...\%'} \\
"מסתיים ב" & \sqlin{LIKE '\%...'} \\
"שם המשפחה שלהם \lr{X}" (בעמודה של שם מלא) & \sqlin{LIKE '\%X\%'} \\
"\textbf{אין לו}" \ \textbar\ \ "לא רשום" \ \textbar\ \ "עדיין פתוח" & \sqlin{IS NULL} \\
"יש לו" \ \textbar\ \ "רשום" \ \textbar\ \ "כבר הסתיים" & \sqlin{IS NOT NULL} \\
"\textbf{וגם}" & \sqlin{AND} \\
"\textbf{או}" & \sqlin{OR} \ \ \textendash{} ואם יש עוד תנאי: \textbf{סוגריים} \\
"לפני \lr{1.7.2026}" & \sqlin{< '2026-07-01'} \\
"אחרי \lr{1.7.2026}" & \sqlin{> '2026-07-01'} \\
\bottomrule
\end{tabular}

\vspace{7pt}
\needspace{9\baselineskip}
\noindent\textbf{\color{inkblue}ג. הסדר \textendash{} באיזה סדר להציג}\par\vspace{2pt}
\noindent\begin{tabular}{@{}p{0.46\linewidth}p{0.48\linewidth}@{}}
\toprule
\thd{מה כתוב בשאלה} & \thd{מה כותבים} \\
\midrule
"ממוינים לפי המחיר" & \sqlin{ORDER BY BookPrice} \\
"מהיקר לזול" \ \textbar\ \ "מהגדול לקטן" \ \textbar\ \ "יורד" & \sqlin{DESC} \\
"מהזול ליקר" \ \textbar\ \ "לפי האלף-בית" \ \textbar\ \ "עולה" & בלי כלום \textendash{} \lr{\texttt{ASC}} הוא ברירת המחדל \\
"לפי \lr{X}, ובתוך כל \lr{X} לפי \lr{Y}" & \sqlin{ORDER BY X, Y} \\
"שלושת ה... ביותר" \ \textbar\ \ "\lr{3} הראשונים" & \sqlin{ORDER BY ... DESC LIMIT 3} \\
\bottomrule
\end{tabular}

\vspace{8pt}
\subsection{שלוש דוגמאות, צעד אחרי צעד}

{\small כל הדוגמאות על \textbf{מועדון הקריאה} שבעמוד \lr{3}.}

\vspace{4pt}
\begin{notepanel}{okgreen}{דוגמה \lr{1} \textendash{} העמודות, התנאי והסדר}
"הציגו את \textbf{שם הספר ואת המחיר}, של הספרים \textbf{שיצאו אחרי \lr{2018}} \textendash{} \textbf{מהיקר לזול}."
\end{notepanel}

\vspace{-2pt}
\noindent\begin{tabular}{@{}p{0.07\linewidth}p{0.40\linewidth}p{0.47\linewidth}@{}}
\toprule
& \thd{מה שואלים את עצמנו} & \thd{מה נכתב} \\
\midrule
\lr{1} & על מה מדובר? ספרים & \sqlin{FROM Books} \\
\lr{2} & מה להציג? שם ומחיר & \sqlin{SELECT BookTitle, BookPrice} \\
\lr{3} & אילו שורות? "ש"יצאו אחרי \lr{2018} & \sqlin{WHERE BookYear > 2018} \\
\lr{4} & באיזה סדר? מהיקר לזול & \sqlin{ORDER BY BookPrice DESC} \\
\bottomrule
\end{tabular}

\begin{otherlanguage}{english}
\begin{sqlgood}
SELECT   BookTitle, BookPrice
FROM     Books
WHERE    BookYear > 2018
ORDER BY BookPrice DESC;
\end{sqlgood}
\end{otherlanguage}

\begin{outplain}
\begin{tabular}{ll}
\toprule
\thd{BookTitle} & \thd{BookPrice} \\
\midrule
ספר המתכונים & 120 \\
סודות הגלקסיה & 95 \\
הדרך לירכא & 72 \\
שירים מהגליל & 45 \\
\bottomrule
\end{tabular}
\end{outplain}

{\small\textbf{שימו לב:} "אחרי \lr{2018}" פירושו \sqlin{> 2018}, ולכן שלושת הספרים משנת \lr{2018} \textbf{אינם} בפלט. אילו היה כתוב "משנת \lr{2018} ואילך" היינו כותבים \sqlin{>= 2018}.}

\vspace{6pt}
\begin{notepanel}{okgreen}{דוגמה \lr{2} \textendash{} \lr{\texttt{DISTINCT}} יחד עם תנאי}
"\textbf{אילו מדפים} יש לספרים \textbf{שמחירם מתחת ל-\lr{50}}? כל מדף \textbf{פעם אחת}."
\end{notepanel}

\vspace{-2pt}
\textbf{הדרך לחשוב על זה:} קודם \textbf{מסננים} (\lr{\texttt{WHERE}}), ורק מה שנשאר עובר ל-\lr{\texttt{DISTINCT}} \textendash{} בדיוק כמו סדר הביצוע מפרק \lr{2}. שלוש שורות עומדות בתנאי:

\begin{outtable}{ של \sqlin{SELECT BookNum, BookShelf, BookPrice FROM Books WHERE BookPrice < 50}}
\begin{tabular}{lll}
\toprule
\thd{BookNum} & \thd{BookShelf} & \thd{BookPrice} \\
\midrule
3 & \nullv & 45 \\
6 & B & 45 \\
7 & A & 38 \\
\bottomrule
\end{tabular}
\end{outtable}

\begin{otherlanguage}{english}
\begin{sqlgood}
SELECT DISTINCT BookShelf
FROM   Books
WHERE  BookPrice < 50;
\end{sqlgood}
\end{otherlanguage}

\begin{outplain}
\begin{tabular}{l}
\toprule
\thd{BookShelf} \\
\midrule
\nullv \\
B \\
A \\
\bottomrule
\end{tabular}
\end{outplain}

{\small\textbf{שתי נקודות:} \lr{(1)} ה-\lr{\texttt{NULL}} \textbf{מופיע בפלט} \textendash{} וזה נכון, כי "אין מדף" הוא אחד המצבים הקיימים בעמודה. \lr{(2)} הפלט \textbf{אינו ממוין} \textendash{} \lr{\texttt{DISTINCT}} מסלק כפילויות בלבד.}

\vspace{6pt}
\begin{notepanel}{okgreen}{דוגמה \lr{3} \textendash{} חישוב, כינוי, ו"שלושת ה... ביותר"}
"הציגו את \textbf{שם הספר} ואת \textbf{המחיר אחרי הנחה של \lr{20\%}}, בכותרת \lr{\texttt{SalePrice}}, עבור \textbf{שלושת הספרים היקרים ביותר}."
\end{notepanel}

\vspace{-2pt}
\noindent\begin{tabular}{@{}p{0.07\linewidth}p{0.40\linewidth}p{0.47\linewidth}@{}}
\toprule
& \thd{מה שואלים את עצמנו} & \thd{מה נכתב} \\
\midrule
\lr{1} & על מה? ספרים & \sqlin{FROM Books} \\
\lr{2} & הנחה של \lr{20\%}, כלומר נשאר \lr{80\%} & \sqlin{BookPrice * 0.8 AS SalePrice} \\
\lr{3} & תנאי? \textbf{אין} \textendash{} "שלושת היקרים" הוא סדר, לא תנאי & בלי \lr{\texttt{WHERE}} \\
\lr{4} & "היקרים ביותר" \lr{$\rightarrow$} יורד, ואז לקצוץ ל-\lr{3} & \lr{\texttt{ORDER BY BookPrice DESC} \ + \ \texttt{LIMIT 3}} \\
\bottomrule
\end{tabular}

\begin{otherlanguage}{english}
\begin{sqlgood}
SELECT   BookTitle, BookPrice * 0.8 AS SalePrice
FROM     Books
ORDER BY BookPrice DESC
LIMIT    3;
\end{sqlgood}
\end{otherlanguage}

\begin{outplain}
\begin{tabular}{ll}
\toprule
\thd{BookTitle} & \thd{SalePrice} \\
\midrule
ספר המתכונים & 96.0 \\
סודות הגלקסיה & 76.0 \\
הדרך לירכא & 57.6 \\
\bottomrule
\end{tabular}
\end{outplain}

\begin{notepanel}{warnred}{המלכודת של דוגמה \lr{3}}
"שלושת הספרים היקרים ביותר" נשמע כמו תנאי, ולכן הרבה תלמידים מנסים לכתוב \lr{\texttt{WHERE}}. \textbf{אין שום תנאי בשאלה הזאת} \textendash{} אי אפשר לדעת מראש מאיזה מחיר מתחילים "היקרים". \textbf{ממיינים, ואז קוצצים.} זה תמיד \lr{\texttt{ORDER BY} \ + \ \texttt{LIMIT}}.

\vspace{3pt}
וגם: \textbf{הנחה של \lr{20\%} היא הכפלה ב-\lr{0.8}}, לא ב-\lr{0.2}. הכפלה ב-\lr{0.2} נותנת את \textbf{גודל ההנחה}, ולא את המחיר לתשלום.
\end{notepanel}

\newpage
\section{דף העזר \textendash{} הכול בטבלה אחת}
\rulebar

{\small זה הפרק שתסתכלו בו בזמן הבחינה. הדוגמאות על בסיס הנתונים של מועדון הקריאה.}

\vspace{4pt}
\noindent\begin{tabular}{@{}p{0.235\linewidth}p{0.46\linewidth}p{0.24\linewidth}@{}}
\toprule
\thd{מה רוצים} & \thd{איך כותבים} & \thd{הערה} \\
\midrule
\multicolumn{3}{@{}l}{\textbf{\lr{SELECT} \textendash{} אילו עמודות}}\\
\midrule
כל העמודות & \sqlin{SELECT * FROM Books} & מסוכן בבחינה \\
עמודות נבחרות & \sqlin{SELECT BookTitle, BookPrice FROM Books} & הסדר שלכם \\
כותרת אחרת & \sqlin{SELECT BookPrice AS Price} & מילה אחת \\
כותרת בעברית & \sqlin{SELECT BookPrice AS "מחיר"} & גרש \textbf{כפול} \\
חישוב & \sqlin{SELECT BookPrice * 1.18 AS WithVat} & לכל שורה \\
חיבור טקסט & \sqlin{SELECT ReaderFirst || ' ' || ReaderLast} & רווח באמצע \\
בלי כפילויות & \sqlin{SELECT DISTINCT BookYear FROM Books} & על כל השורה \\
\midrule
\multicolumn{3}{@{}l}{\textbf{\lr{WHERE} \textendash{} אילו שורות}}\\
\midrule
שווה / לא שווה & \lr{\texttt{WHERE BookYear = 2018} \ \textbar\ \ \texttt{<> 2018}} & \sqlin{=} אחד \\
גדול / קטן & \sqlin{WHERE BookPrice >= 70} & \\
טקסט & \sqlin{WHERE BookShelf = 'A'} & גרש \textbf{בודד} \\
תאריך & \sqlin{WHERE JoinDate >= '2025-01-01'} & \lr{\texttt{YYYY-MM-DD}} \\
טווח & \sqlin{WHERE BookPrice BETWEEN 45 AND 72} & \textbf{כולל} קצוות \\
רשימה & \sqlin{WHERE BookYear IN (2019, 2023)} & קיצור ל-\lr{\texttt{OR}} \\
לא ברשימה & \sqlin{WHERE BookYear NOT IN (2018)} & \\
מתחיל ב & \sqlin{WHERE BookTitle LIKE 'ס\%'} & \\
מכיל & \sqlin{WHERE BookTitle LIKE '\%גל\%'} & \textbf{שני} אחוזים \\
תו אחד & \sqlin{WHERE BookShelf LIKE '\_'} & \\
ערך חסר & \sqlin{WHERE BookShelf IS NULL} & \textbf{לא} \sqlin{= NULL} \\
ערך קיים & \sqlin{WHERE ReaderPhone IS NOT NULL} & \\
שני תנאים & \sqlin{WHERE BookYear = 2018 AND BookPrice > 60} & מצמצם \\
אחד מהשניים & \sqlin{WHERE BookShelf = 'A' OR BookShelf = 'C'} & מרחיב \\
\textbf{\lr{OR} עם עוד תנאי} & \sqlin{WHERE (A OR B) AND C} & \textbf{סוגריים!} \\
\midrule
\multicolumn{3}{@{}l}{\textbf{\lr{ORDER BY} \textendash{} באיזה סדר}}\\
\midrule
עולה & \sqlin{ORDER BY BookPrice} & ברירת מחדל \\
יורד & \sqlin{ORDER BY BookPrice DESC} & חובה לכתוב \\
שני מפתחות & \sqlin{ORDER BY BookYear DESC, BookPrice} & \lr{\texttt{DESC}} לכל אחד \\
לפי כינוי & \sqlin{ORDER BY AfterSale DESC} & מותר כאן \\
לפי מספר עמודה & \sqlin{ORDER BY 2 DESC} & שביר \\
\lr{N} הראשונים & \sqlin{ORDER BY BookPrice DESC LIMIT 3} & תמיד עם מיון \\
\bottomrule
\end{tabular}

\vspace{7pt}
\begin{panel}
\textbf{השלד שאפשר לכתוב מההתחלה, בכל שאלה:}

\vspace{3pt}
\begin{otherlanguage}{english}
\begin{sqlverb}
SELECT   
FROM     
WHERE    
ORDER BY ;
\end{sqlverb}
\end{otherlanguage}
\vspace{1pt}
ממלאים מלמעלה למטה, ומוחקים את מה שלא צריך. \textbf{ונקודה-פסיק בסוף.}
\end{panel}

\newpage
\section{הטעויות שחוזרות על עצמן}
\rulebar

\vspace{2pt}
\noindent\begin{tabular}{@{}p{0.04\linewidth}p{0.29\linewidth}p{0.30\linewidth}p{0.30\linewidth}@{}}
\toprule
& \thd{הטעות} & \thd{מה קורה בפועל} & \thd{התיקון} \\
\midrule
\lr{1} & \sqlin{WHERE X = NULL} & \textbf{אפס שורות}, בלי שגיאה & \sqlin{WHERE X IS NULL} \\
\lr{2} & גרש כפול למחרוזת: \sqlin{= "A"} & מחפש \textbf{עמודה} בשם \lr{A} & גרש בודד: \sqlin{= 'A'} \\
\lr{3} & \lr{\texttt{OR}} בלי סוגריים & \lr{\texttt{AND}} גובר, ונכנסות שורות מיותרות & \sqlin{(A OR B) AND C} \\
\lr{4} & \sqlin{LIKE 'גליל'} & התאמה מדויקת בלבד & \sqlin{LIKE '\%גליל\%'} \\
\lr{5} & \lr{\texttt{ORDER BY}} לפני \lr{\texttt{WHERE}} & \textbf{שגיאת תחביר}, לא רץ & הסדר: \lr{\texttt{SELECT-FROM-WHERE-ORDER BY}} \\
\lr{6} & \lr{\texttt{DESC}} על העמודה הראשונה בלבד & העמודה השנייה נשארה עולה & \lr{\texttt{DESC}} ליד \textbf{כל} עמודה שצריך \\
\lr{7} & פסיק לפני \lr{\texttt{FROM}} & \lr{\texttt{syntax error near FROM}} & מוחקים את הפסיק \\
\lr{8} & סינון לפי כינוי ב-\lr{\texttt{WHERE}} & \lr{\texttt{no such column}} & כותבים את החישוב במלואו \\
\lr{9} & \sqlin{BETWEEN 72 AND 45} & אפס שורות, בלי שגיאה & הקטן \textbf{קודם} \\
\lr{10} & \sqlin{SELECT *} כשביקשו שתי עמודות & מציג את כל העמודות & בדיוק מה שביקשו \\
\lr{11} & אין נקודה-פסיק בסוף & השאילתה "נתקעת" ומחכה להמשך & \sqlin{;} בסוף כל שאילתה \\
\bottomrule
\end{tabular}

\vspace{8pt}
\begin{panel}
\textbf{נתקעתם בשגיאה?} קראו את המילה שאחרי \lr{\texttt{near}} \textendash{} הבעיה נמצאת שם, או ממש לפניה. \lr{\texttt{no such column: X}} פירושו כמעט תמיד שכתבתם מחרוזת בגרש כפול, או שטעיתם באיות שם העמודה.
\end{panel}

\vspace{6pt}
\section{תרגול א' \textendash{} כתיבת שאילתות}
\rulebar

{\small על \textbf{מועדון הקריאה} שבעמוד \lr{3}. כסו את התשובות, פתרו על נייר, ואז בדקו.}

\vspace{3pt}
\subq{1}{הציגו את שם הספר ואת המחיר, של כל הספרים שמחירם בין \lr{45} ל-\lr{70} \textendash{} מהיקר לזול.}
\subq{2}{הציגו את מספרי הספרים שאין להם מדף.}
\subq{3}{אילו מדפים בכלל קיימים במועדון? כל מדף פעם אחת.}
\subq{4}{הציגו שם מלא של כל קורא (שם פרטי, רווח, שם משפחה), בכותרת \lr{\texttt{FullName}}.}
\subq{5}{הציגו את שמות הספרים שמופיעות בכותרתם האותיות "גל", ושיצאו \textbf{לפני} שנת \lr{2020}.}
\subq{6}{הציגו את שם הספר ואת המחיר אחרי מע"מ (\lr{18\%}), בכותרת "מחיר סופי", לספרים ממדף \lr{A} או \lr{B} שמחירם מעל \lr{60}.}
\subq{7}{הציגו את שלושת הספרים הזולים ביותר.}
\subq{8}{הציגו את הקוראים שאין להם טלפון, לפי תאריך הצטרפות מהמוקדם למאוחר.}

\vspace{8pt}
\begin{panel}
\textbf{התשובות}

\vspace{2pt}
\begin{otherlanguage}{english}
\begin{Verbatim}[fontsize=\small]
1.  SELECT BookTitle, BookPrice FROM Books
    WHERE BookPrice BETWEEN 45 AND 70 ORDER BY BookPrice DESC;

2.  SELECT BookNum FROM Books WHERE BookShelf IS NULL;

3.  SELECT DISTINCT BookShelf FROM Books;

4.  SELECT ReaderFirst || ' ' || ReaderLast AS FullName FROM Readers;

5.  SELECT BookTitle FROM Books WHERE BookTitle LIKE '%גל%' AND BookYear < 2020;

6.  SELECT BookTitle, BookPrice * 1.18 AS "מחיר סופי" FROM Books
    WHERE (BookShelf = 'A' OR BookShelf = 'B') AND BookPrice > 60;

7.  SELECT BookTitle, BookPrice FROM Books ORDER BY BookPrice LIMIT 3;

8.  SELECT ReaderFirst, ReaderLast FROM Readers
    WHERE ReaderPhone IS NULL ORDER BY JoinDate;
\end{Verbatim}
\end{otherlanguage}

\vspace{2pt}
{\small\textbf{שימו לב:} בסעיף \lr{2} \textbf{אסור} \sqlin{= NULL}. בסעיף \lr{3} המדף החסר יופיע גם הוא, כשורת \lr{\texttt{NULL}} \textendash{} זה נכון, כי זה אחד הערכים הקיימים בעמודה. בסעיף \lr{6} \textbf{הסוגריים הכרחיים}. בסעיף \lr{7} "הזולים" פירושו \lr{\texttt{ASC}}, כלומר בלי \lr{\texttt{DESC}}.}
\end{panel}

\newpage
\section{תרגול מסוג הבחינה}
\rulebar

התרגול שבפרק הקודם בדק דבר אחד: \textbf{לכתוב שאילתה}. אבל בבחינה יש עוד \textbf{שלושה סוגי שאלות}, והם שליש מהנקודות: שאלות על \textbf{מבנה} בסיס הנתונים, \textbf{תיקון} שאילתה שגויה, ו\textbf{ניתוח פלט}. הפרק הזה מתרגל בדיוק אותם.

\vspace{3pt}
\begin{panel}
\textbf{שימו לב לתנאי התרגול.} חלק א' רץ על בסיס נתונים \textbf{שלא ראיתם} \textendash{} וזה בכוונה. בבחינה תקבלו סכמה חדשה לגמרי, וצריך להתאמן על לקרוא אותה מאפס, ולא על לזכור אותה.
\end{panel}

\subsection{חלק א' \textendash{} מבנה בסיס הנתונים}

\textbf{בית הספר לנהיגה "דרך בטוחה".} לפניכם חמש טבלאות. לכל תלמיד רשומה עיר מגורים אחת; לכל מורה רכב אחד, ולכל רכב מורה אחד; ותלמיד יכול ללמוד אצל כמה מורים, וכל מורה מלמד כמה תלמידים.

\vspace{3pt}
\begin{otherlanguage}{english}
\begin{center}
\begin{tikzpicture}[font=\scriptsize\sffamily]
\tikzset{tblm/.style={matrix of nodes,
  nodes={anchor=west, inner sep=2.6pt, text width=24mm, align=left, font=\scriptsize\sffamily},
  row sep=-\pgflinewidth, draw=inkblue, line width=0.7pt, inner sep=0pt,
  row 1/.style={nodes={fill=inkblue, text=white, font=\scriptsize\sffamily\bfseries,
                       align=center, text width=24mm}}}}
\matrix (T) [tblm] { Towns \\ TownCode \\ TownName \\ };
\matrix (P) [tblm, right=9mm of T] { Pupils \\ PupilNum \\ PupilName \\ PupilTown \\ BirthYear \\ };
\matrix (L) [tblm, right=9mm of P] { PupilTeach \\ LinkPupil \\ LinkTeach \\ Lessons \\ };
\matrix (E) [tblm, right=9mm of L] { Teachers \\ TeachNum \\ TeachName \\ TeachCar \\ Seniority \\ };
\matrix (C) [tblm, right=9mm of E] { Cars \\ CarCode \\ CarModel \\ CarGear \\ };
\end{tikzpicture}
\end{center}
\end{otherlanguage}

\vspace{2pt}
{\small\textbf{דוגמה לשורות בטבלה \lr{\texttt{PupilTeach}}:} \lr{(201, 7, 12)}, \lr{(201, 9, 4)}, \lr{(205, 7, 20)} \textendash{} כלומר תלמיד \lr{201} למד \lr{12} שיעורים אצל מורה \lr{7}, ועוד \lr{4} שיעורים אצל מורה \lr{9}.}

\vspace{3pt}
\subq{1}{כתבו את \textbf{המפתח הראשי} של חמש הטבלאות. באחת מהן המפתח \textbf{מורכב}.}
\subq{2}{כתבו את \textbf{כל המפתחות הזרים} בבסיס הנתונים, ולאיזה שדה כל אחד מצביע.}
\subq{3}{מהו \textbf{סוג הקשר} בכל זוג: \lr{\texttt{Towns}\,\textendash\,\texttt{Pupils} \ \textbar\ \ \texttt{Pupils}\,\textendash\,\texttt{Teachers} \ \textbar\ \ \texttt{Teachers}\,\textendash\,\texttt{Cars}}? נמקו כל אחד במשפט.}
\subq{4}{נכון או לא נכון: "יכולה להיות עיר בטבלת \lr{\texttt{Towns}} שאף תלמיד אינו גר בה". הסבירו.}

\subsection{חלק ב' \textendash{} תיקון שאילתה שגויה}

{\small על \textbf{מועדון הקריאה} שבעמוד \lr{3}.}

\vspace{2pt}
\subq{5}{בשאילתה הבאה \textbf{שלוש} טעויות. כתבו אותה מתוקנת. המטרה: שם ומחיר של הספרים שמחירם מעל \lr{50}, מהיקר לזול.}

\begin{otherlanguage}{english}
\begin{sqlbad}
SELECT BookTitle, BookPrice,
FROM   Books
ORDER BY BookPrice DESC
WHERE  BookPrice > 50
\end{sqlbad}
\end{otherlanguage}

\subq{6}{בשאילתה הבאה \textbf{שתי} טעויות. המטרה: הקוראים שאין להם טלפון, מהמצטרף האחרון לראשון.}

\begin{otherlanguage}{english}
\begin{sqlbad}
SELECT   ReaderFirst, ReaderLast
FROM     Readers
WHERE    ReaderPhone = NULL
ORDER BY JoinDate DESC
\end{sqlbad}
\end{otherlanguage}

\subsection{חלק ג' \textendash{} ניתוח פלט}

\subq{7}{לפניכם שאילתה. \lr{(א)} כמה שורות היא תחזיר? \lr{(ב)} איזה ספר יופיע \textbf{אחרון}, ומדוע?}

\begin{otherlanguage}{english}
\begin{sqlverb}
SELECT   BookTitle, BookShelf
FROM     Books
ORDER BY BookShelf DESC;
\end{sqlverb}
\end{otherlanguage}

\subq{8}{לפניכם \textbf{הפלט} של שאילתה. כתבו את השאילתה הפשוטה ביותר שמפיקה אותו.}

\begin{outplain}
\begin{tabular}{l}
\toprule
\thd{BookShelf} \\
\midrule
A \\
B \\
\nullv \\
C \\
\bottomrule
\end{tabular}
\end{outplain}

\newpage
\subsection{התשובות לחלק א'\textendash ג'}

\begin{panel}
\textbf{\lr{1.} מפתחות ראשיים}

\vspace{2pt}
\noindent\begin{tabular}{@{}p{0.24\linewidth}p{0.70\linewidth}@{}}
\lr{\texttt{Towns}} & \lr{\texttt{TownCode}} \\
\lr{\texttt{Pupils}} & \lr{\texttt{PupilNum}} \\
\lr{\texttt{Teachers}} & \lr{\texttt{TeachNum}} \\
\lr{\texttt{Cars}} & \lr{\texttt{CarCode}} \\
\lr{\texttt{PupilTeach}} & \lr{\texttt{LinkPupil} + \texttt{LinkTeach}} \textendash{} \textbf{מפתח מורכב} \\
\end{tabular}

\vspace{3pt}
{\small\textbf{למה מורכב?} בדוגמה, \lr{201} מופיע פעמיים ו-\lr{7} מופיע פעמיים \textendash{} אף אחד מהם לבדו אינו ייחודי. רק \textbf{הצירוף} מופיע פעם אחת.}
\end{panel}

\begin{panel}
\textbf{\lr{2.} מפתחות זרים}

\vspace{2pt}
\noindent\begin{tabular}{@{}p{0.30\linewidth}p{0.30\linewidth}p{0.34\linewidth}@{}}
\toprule
\thd{המפתח הזר} & \thd{בטבלה} & \thd{מצביע אל} \\
\midrule
\lr{\texttt{PupilTown}} & \lr{\texttt{Pupils}} & \lr{\texttt{TownCode}} \\
\lr{\texttt{TeachCar}} & \lr{\texttt{Teachers}} & \lr{\texttt{CarCode}} \\
\lr{\texttt{LinkPupil}} & \lr{\texttt{PupilTeach}} & \lr{\texttt{PupilNum}} \\
\lr{\texttt{LinkTeach}} & \lr{\texttt{PupilTeach}} & \lr{\texttt{TeachNum}} \\
\bottomrule
\end{tabular}

\vspace{3pt}
{\small בטבלה מקשרת, \textbf{שני} השדות הם מפתחות זרים \textendash{} וביחד הם גם המפתח הראשי. \lr{\texttt{Lessons}} אינו מפתח זר: הוא סתם מספר שיעורים.}
\end{panel}

\begin{panel}
\textbf{\lr{3.} סוגי הקשרים}

\vspace{2pt}
\noindent\begin{tabular}{@{}p{0.26\linewidth}p{0.10\linewidth}p{0.58\linewidth}@{}}
\toprule
\thd{הזוג} & \thd{הקשר} & \thd{הנימוק} \\
\midrule
\lr{\texttt{Towns}\,\textendash\,\texttt{Pupils}} & \textbf{\lr{1:N}} & בעיר אחת גרים כמה תלמידים, ולכל תלמיד עיר אחת. המפתח הזר \lr{\texttt{PupilTown}} יושב בצד הרבים \\
\lr{\texttt{Pupils}\,\textendash\,\texttt{Teachers}} & \textbf{\lr{M:N}} & תלמיד לומד אצל כמה מורים, ומורה מלמד כמה תלמידים. לכן יש טבלה מקשרת, \lr{\texttt{PupilTeach}} \\
\lr{\texttt{Teachers}\,\textendash\,\texttt{Cars}} & \textbf{\lr{1:1}} & לכל מורה רכב אחד, ולכל רכב מורה אחד. \lr{\texttt{TeachCar}} הוא ייחודי \\
\bottomrule
\end{tabular}

\vspace{3pt}
{\small\textbf{החצים ב-\lr{DSD}:} \lr{\texttt{PupilTown} $\rightarrow$ \texttt{TownCode}} \ \textbullet\ \ \lr{\texttt{TeachCar} $\rightarrow$ \texttt{CarCode}} \ \textbullet\ \ \lr{\texttt{LinkPupil} $\rightarrow$ \texttt{PupilNum}} \ \textbullet\ \ \lr{\texttt{LinkTeach} $\rightarrow$ \texttt{TeachNum}}}
\end{panel}

\begin{panel}
\textbf{\lr{4.} נכון.} בקשר \lr{1:N} מותר לצד ה"אחד" להישאר בלי אף שורה מתאימה: עיר שרשומה במערכת, ואף תלמיד אינו גר בה. \textbf{ההפך אסור} \textendash{} אי אפשר לרשום תלמיד עם קוד עיר שאינו קיים בטבלת \lr{\texttt{Towns}}.
\end{panel}

\begin{panel}
\textbf{\lr{5.} שלוש הטעויות:} \lr{(1)} פסיק מיותר אחרי \lr{\texttt{BookPrice}}, לפני \lr{\texttt{FROM}}. \lr{(2)} \lr{\texttt{ORDER BY}} נכתב \textbf{לפני} \lr{\texttt{WHERE}} \textendash{} הסדר חייב להיות \lr{\texttt{WHERE}} ואז \lr{\texttt{ORDER BY}}. \lr{(3)} חסרה נקודה-פסיק.

\begin{otherlanguage}{english}
\begin{sqlgood}
SELECT   BookTitle, BookPrice
FROM     Books
WHERE    BookPrice > 50
ORDER BY BookPrice DESC;
\end{sqlgood}
\end{otherlanguage}

\begin{outplain}
\begin{tabular}{ll}
\toprule
\thd{BookTitle} & \thd{BookPrice} \\
\midrule
ספר המתכונים & 120 \\
סודות הגלקסיה & 95 \\
הדרך לירכא & 72 \\
מסע אל הירח & 68 \\
\bottomrule
\end{tabular}
\end{outplain}
\end{panel}

\begin{panel}
\textbf{\lr{6.} שתי הטעויות:} \lr{(1)} \sqlin{= NULL} במקום \sqlin{IS NULL} \textendash{} כך נכתב, הפלט ריק. \lr{(2)} חסרה נקודה-פסיק.

\begin{otherlanguage}{english}
\begin{sqlgood}
SELECT   ReaderFirst, ReaderLast
FROM     Readers
WHERE    ReaderPhone IS NULL
ORDER BY JoinDate DESC;
\end{sqlgood}
\end{otherlanguage}

\begin{outplain}
\begin{tabular}{lll}
\toprule
\thd{ReaderFirst} & \thd{ReaderLast} & \thd{(JoinDate)} \\
\midrule
איתי & ברק & 2025-02-10 \\
עומר & לוי & 2023-11-15 \\
\bottomrule
\end{tabular}
\end{outplain}

{\small\lr{\texttt{JoinDate}} מופיע כאן רק כדי להראות את המיון \textendash{} בשאילתה עצמה לא ביקשו אותו, ולכן הוא אינו ב-\lr{\texttt{SELECT}}.}
\end{panel}

\begin{panel}
\textbf{\lr{7.} \lr{(א)} \ \lr{7} שורות} \textendash{} כל הספרים. אין \lr{\texttt{WHERE}}, ו-\lr{\texttt{ORDER BY}} אינו מסנן.

\vspace{2pt}
\textbf{\lr{(ב)} אחרון יופיע "יומן מהמחסן"} (ספר \lr{3}) \textendash{} הספר שאין לו מדף. במיון \textbf{יורד} ערכי \lr{\texttt{NULL}} נופלים \textbf{בסוף}.
\end{panel}

\begin{panel}
\textbf{\lr{8.} השאילתה:}

\begin{otherlanguage}{english}
\begin{sqlgood}
SELECT DISTINCT BookShelf FROM Books;
\end{sqlgood}
\end{otherlanguage}

\vspace{2pt}
{\small שלושה רמזים שמובילים לתשובה: \lr{(1)} \textbf{עמודה אחת} בפלט \lr{$\rightarrow$} \lr{\texttt{SELECT}} של שדה אחד. \lr{(2)} כל מדף מופיע \textbf{פעם אחת}, למרות שבטבלה יש \lr{7} ספרים \lr{$\rightarrow$} \lr{\texttt{DISTINCT}}. \lr{(3)} הפלט \textbf{אינו ממוין} (\lr{A}, \lr{B}, \lr{\textit{NULL}}, \lr{C}) \lr{$\rightarrow$} אין \lr{\texttt{ORDER BY}}. ושורת ה-\lr{\texttt{NULL}} מוכיחה שלא היה \lr{\texttt{WHERE}} שסינן אותה.}
\end{panel}

\vspace{8pt}
{\color{inkblue}\hrule height 1pt}
\vspace{4pt}
\begin{center}
{\small\textbf{בהצלחה בבחינה.} מי שפתר את שני התרגולים \textendash{} \lr{16} השאלות \textendash{} בלי להציץ, מוכן לכל ארבע השאלות שבבחינה.}
\end{center}

\end{document}
